\documentclass{article}

\usepackage[a4paper]{geometry}
\usepackage[english]{babel}
\usepackage{parskip}

\usepackage{amsmath, amssymb}
\usepackage{graphicx}
\usepackage{subcaption}
\usepackage{float}
\usepackage[hidelinks]{hyperref}

\usepackage{listings}
\usepackage{xcolor}

\definecolor{codegreen}{rgb}{0,0.6,0}
\definecolor{codegray}{rgb}{0.5,0.5,0.5}
\definecolor{codepurple}{rgb}{0.58,0,0.82}
\definecolor{backcolour}{rgb}{0.95,0.95,0.92}

\lstdefinestyle{mystyle}{
  backgroundcolor=\color{backcolour},
  commentstyle=\color{codegreen},
  keywordstyle=\color{magenta},
  numberstyle=\tiny\color{codegray},
  stringstyle=\color{codepurple},
  basicstyle=\ttfamily\footnotesize,
  breaklines=true,
  captionpos=b,
  keepspaces=true,
  numbers=left,
  numbersep=5pt,
  showstringspaces=false,
  tabsize=2
}
\lstset{style=mystyle}

\newcommand{\CodeDir}{Code}
\newcommand{\PlotDir}{Plots}
\newcommand{\CalcDir}{Calculations}

\title{NURa Hand-in 3}
\author{Silvan Toet (s1827332)}
\date{\today}

\begin{document}
\maketitle

\section{Satellite galaxies around a massive central}
We once again have the number density profile for satellite galaxies around a massive central (dark matter) halo.
The spherical volume element \(dV\) can be defined as follows:
We know that the volume can be computed in spherical coordinates using the triple integral
\begin{align}
  \iiint_V n(x) dV & = \int_{\phi=0}^{2\pi} \int_{\theta=0}^{\pi} \int_{x=0}^{x_{max}} n(x) x^2 \sin\theta \, dx\,d\theta\,d\phi                                                    \\
                   & = \int_{x=0}^{x_{max}} 4 \pi n(x) x^2 dx = \int_{x=0}^{x_{max}} N(x) dx                                                                                        \\
                   & = 4 \pi A \langle N_{sat} \rangle \int_{x=0}^{x_{max}=5} x^2 \left(\frac{x}{b}\right)^{a-3} \exp(-\left(\frac{x}{b}\right)^c) dx = \langle N_{sat} \rangle     \\
                   & = 4 \pi A \langle N_{sat} \rangle \int_{x=0}^{x_{max}=5} x^{a-1} \left(\frac{1}{b}\right)^{a-3} \exp(-\left(\frac{x}{b}\right)^c) dx = \langle N_{sat} \rangle \\
  -> A             & = \frac{1}{4 \pi \int_{x=0}^{x_{max}=5} x^{a-1} \left(\frac{1}{b}\right)^{a-3} \exp(-\left(\frac{x}{b}\right)^c) dx}
\end{align}
Where we have defined the number of satellites N(x)dx in the infinitesimal range [x, x+dx).
Since n(x) where x is the radius relative to virial radius \(x=r/r_{vir}\) is only dependent on x, we can easily solve for \(\theta\) and \(\phi\).
Dividing by \(\langle N_{sat} \rangle\) on both sides and rearranging gives us an expression for \(A\).

\lstinputlisting[language=Python]{\CodeDir/satellites_general_code.txt}

\subsection{(1a) Maximum of N(x)}
For this question we are asked to find the maximum of N(x), the number of satellites in the infinitesimal range [x, x+dx).
Let us first examine both functions.
\begin{figure}[H]
  \centering
  
\includegraphics[width=0.85\textwidth]{\PlotDir/nx_vs_Nx.png}
  \caption{n(x) dx, the number density profile, compared to N(x) dx, the number of satellites in the infinitesimal range [x, x+dx)}
\end{figure}


Although n(x), the number density profile, is a monotonically decreasing function, because the number density will decrease with radius, N(x) has a maximum because it can first increase at small x while the x^2 term gets larger, before the exponential decay dominates at larger radii.


Explain the maximization algorithm used.

The maximum of \(N(x)\) is found at:
$$x_{\text{max}} \approx  \input{\CalcDir/satellite_max_x.txt} $$
$$N(x_{\text{max}}) \approx  \input{\CalcDir/satellite_max_Nx.txt} $$

\subsubsection{Code for N(x) maximization}
NOTE: all code is reproduced below as an example, show only the relevant lines for this exercise here instead.
\lstinputlisting[language=Python]{\CodeDir/satellites_maximize_code.txt}

\subsection{(1b) $\chi^2$-fit}
For each halo-mass bin, we fit the same parametric profile n(x;a,b,c) separately. This yields five different best-fit parameter sets (a,b,c), one for each mass bin.
The normalization A is not an independent fit parameter, but is recalculated from a,b,c for each model using Romberg integration.
Because the radial coordinate is normalized by the virial radius, the model describes the dimensionless shape of the satellite distribution, allowing that shape to vary with halo mass.

Argue the choices of your binning. Explain the minimization routine used. Bonus points for a correctly implemented derivative-based routine.

\begin{table}[H]
  \centering
  \caption{Best-fit parameters for $\chi^2$ fits.}
  \begin{tabular}{|c|c|c|c|c|c|}
    \hline
    Dataset & $N_{\mathrm{sat}}$ & $\chi^2$ & $a$ & $b$ & $c$ \\
    \hline
    m11 & 0.01368 & 0.00000 & 0.00000 & 0.00000 & 0.00000 \\
m12 & 0.25089 & 0.00000 & 0.00000 & 0.00000 & 0.00000 \\
m13 & 4.37351 & 0.00000 & 0.00000 & 0.00000 & 0.00000 \\
m14 & 29.13453 & 0.00000 & 0.00000 & 0.00000 & 0.00000 \\
m15 & 329.50000 & 0.00000 & 0.00000 & 0.00000 & 0.00000
             \\
    \hline
  \end{tabular}
\end{table}

\begin{figure}[H]
  \centering
  \includegraphics[width=0.85\textwidth]{\PlotDir/satellite_fits_chi2.png}
  \caption{Fit of the satellite number density profile using the chi-squared method.}
\end{figure}

\subsubsection{Code for $\chi^2$-fit}
NOTE: all code is reproduced below as an example, show only the relevant lines for this exercise here instead.
\lstinputlisting[language=Python]{\CodeDir/satellites_chi2_code.txt}

\subsection{(1c) Poisson ln-likelihood fit}

Explain how you determine the Poisson ln-likelihood. Describe the optimization method used. How do your best fit parameters differ from the ones obtained in (1b)? Why?

\begin{table}[H]
  \centering
  \caption{Best-fit parameters for $\chi^2$ fits.}
  \begin{tabular}{|c|c|c|c|c|}
    \hline
    Dataset & llh & $a$ & $b$ & $c$                  \\
    \hline
    \input{Calculations/table_fitparams_poisson.tex} \\
    \hline
  \end{tabular}
\end{table}

\begin{figure}[H]
  \centering
  \includegraphics[width=0.8\textwidth]{\PlotDir/satellite_fits_poisson.png}
  \caption{Fit of the satellite number density profile using the Poisson log-likelihood method.}
\end{figure}

\subsubsection{Code for Poisson ln-likelihood fit}
NOTE: all code is reproduced below as an example, show only the relevant lines for this exercise here instead.
\lstinputlisting[language=Python]{\CodeDir/satellites_poisson_code.txt}

\subsection{(1d) Statistical tests}

% change table to include KS score if you implemented it
\begin{table}[H]
  \centering
  \caption{Statistical tests for the $\chi^2$ and Poisson ln-likelihood fits.}
  \begin{tabular}{|c|c|c|c|}
    \hline
    Method & Dataset & G & Q                         \\
    \hline
    \input{\CalcDir/statistical_test_table_rows.tex} \\
    \hline
  \end{tabular}
\end{table}

What can you conclude from the results of the statistical tests?

\subsubsection{Code for statistical tests}
NOTE: all code is reproduced below as an example, show only the relevant lines for this exercise here instead.
\lstinputlisting[language=Python]{\CodeDir/satellites_statistical_tests_code.txt}

\subsection{(1e) \textbf{Bonus} Monte Carlo simulation}

Pick a single dataset. Describe the implementation of your Monte Carlo parametric bootstrap.

\begin{figure}[H]
  \centering
  \includegraphics[width=0.85\textwidth]{\PlotDir/satellite_monte_carlo_chi2.png}
  \caption{Best $\chi^2$-fit profiles using pseudo-experiments. }
\end{figure}

\begin{figure}[H]
  \centering
  \includegraphics[width=0.85\textwidth]{\PlotDir/satellite_monte_carlo_poisson.png}
  \caption{Best Poisson-likelihood fit profiles using pseudo-experiments. }
\end{figure}

\subsubsection{Code for Monte Carlo simulation}
NOTE: all code is reproduced below as an example, show only the relevant lines for this exercise here instead.
\lstinputlisting[language=Python]{\CodeDir/satellites_monte_carlo_code.txt}

\end{document}
